\chapter{Unterhalbgruppen und allgemeine Konstruktionen}

\section{Erzeugen und Abschließen}

\begin{definition}[Unterhalbgruppe und Erzeugnis]
Eine Teilmenge $A\subseteq S$ heißt \emph{Unterhalbgruppe}, wenn
$AA\subseteq A$ gilt, wobei
\[
 AB=\{ab:a\in A,\ b\in B\}\qquad(A,B\subseteq S).
\]
Die Multiplikation auf $A$ ist die Einschränkung derjenigen von $S$.
Für $X\subseteq S$ bezeichnet $\langle X\rangle$ die kleinste
Unterhalbgruppe, die $X$ enthält. Eine solche Menge $X$ heißt
\emph{Erzeugendensystem} von $\langle X\rangle$.
\end{definition}

\begin{theorem}[Existenz und Beschreibung des Erzeugnisses]
\[
 \langle X\rangle
 =\bigcap_{\substack{A\subseteq S\\X\subseteq A,\ AA\subseteq A}} A
 =\bigcup_{n\geq1}X^n ,
 \qquad X^n=\{x_1\cdots x_n:x_i\in X\}.
\]
Insbesondere ist $\langle\varnothing\rangle=\varnothing$.
Für $X\subseteq Y\subseteq S$ gelten
$X\subseteq\langle X\rangle$,
$\langle X\rangle\subseteq\langle Y\rangle$ und
$\langle\langle X\rangle\rangle=\langle X\rangle$.
\cite{harju}
\end{theorem}

\begin{theorem}[Durchschnitte und gerichtete Vereinigungen]
Jeder Durchschnitt von Unterhalbgruppen von $S$ ist eine
Unterhalbgruppe; der Durchschnitt der leeren Familie ist $S$.
Ist $(A_i)_{i\in I}$ eine nichtleere Familie von Unterhalbgruppen
mit der Eigenschaft, daß es zu $i,j\in I$ ein $k\in I$ mit
$A_i\cup A_j\subseteq A_k$ gibt, so ist auch
$\bigcup_{i\in I}A_i$ eine Unterhalbgruppe.
Die Unterhalbgruppen bilden unter Inklusion einen vollständigen Verband,
mit
\[
 \bigwedge_i A_i=\bigcap_i A_i,\qquad
 \bigvee_i A_i=\left\langle\bigcup_i A_i\right\rangle .
\]
\cite{howie}
\end{theorem}

\begin{definition}[Kompakte Elemente und algebraische Verbände]
\label{b28guide:def:compact}
In einem vollständigen Verband $L$ heißt $c\in L$ \emph{kompakt},
wenn für jede Teilmenge $M\subseteq L$ aus $c\leq\bigvee M$ bereits
$c\leq\bigvee F$ für eine endliche Teilmenge $F\subseteq M$ folgt.
Der Verband heißt \emph{algebraisch}, wenn jedes seiner Elemente
ein Supremum kompakter Elemente ist.
\end{definition}

\begin{theorem}[Algebraischer Unterhalbgruppenverband]
\label{b28guide:thm:algebraic-subsemigroups}
Für jede Halbgruppe $S$ ist der Verband ihrer Unterhalbgruppen
algebraisch. Seine kompakten Elemente sind genau die endlich erzeugten
Unterhalbgruppen. Für jede Teilmenge $X\subseteq S$ gilt
\[
 \langle X\rangle
 =\bigcup_{\substack{F\subseteq X\\F\text{ endlich}}}\langle F\rangle.
\]
Insbesondere ist $S$ die gerichtete Vereinigung seiner endlich erzeugten
Unterhalbgruppen. Endliche Erzeugung bedeutet hier Erzeugung durch
endlich viele Elemente; weder $S$ noch diese Unterhalbgruppen
müssen endlich sein. \cite{burris}
\end{theorem}

\section{Produkte, Umkehrung und Erweiterungen}

\begin{definition}[Direktes Produkt und Gegenhalbgruppe]
Für eine mengenindizierte Familie $(S_i)_{i\in I}$ ist das
\emph{direkte Produkt} $\prod_{i\in I}S_i$ mit
$(s_i)_i(t_i)_i=(s_it_i)_i$ versehen.
Das Produkt der leeren Familie ist die einelementige Halbgruppe.
Die \emph{Gegenhalbgruppe} $S^{\mathrm{op}}$ hat die Menge $S$ und
die Multiplikation $a\circ b=ba$.
\end{definition}

\begin{theorem}[Abschluß unter Produkten und Umkehrung]
Jedes direkte Produkt von Halbgruppen und jede Gegenhalbgruppe
ist wieder eine Halbgruppe. Es gilt
$(S^{\mathrm{op}})^{\mathrm{op}}=S$.
Das Umkehren aller Produkte in einer allgemeingültigen
Halbgruppenaussage liefert wieder eine allgemeingültige Aussage;
dabei werden insbesondere Links- und Rechtsbegriffe vertauscht.
\cite{harju,howie}
\end{theorem}

\begin{theorem}[Adjunktion einer Einheit oder eines Nullelements]
Die Erweiterung $S^1$ aus Definition~\ref{b28guide:def:unit}
ist eine Halbgruppe und enthält $S$ als Unterhalbgruppe.
Ebenso ergibt die Adjunktion eines neuen Elements $0$ zu $S$
mit $0x=x0=0$ eine Halbgruppe $S^0$.
Beide Konstruktionen lassen die Produkte von Elementen aus $S$ unverändert.
\cite{harju,gould}
\end{theorem}

\section{Die Potenzhalbgruppe}

\begin{definition}[Potenzhalbgruppe]
Die Potenzmenge $\mathcal P(S)$ wird mit dem Teilmengenprodukt
$(A,B)\mapsto AB$ versehen.
$\mathcal P_+(S)=\mathcal P(S)\setminus\{\varnothing\}$
bezeichnet die Menge der nichtleeren Teilmengen.
\end{definition}

\begin{theorem}[Assoziativität und Einbettung durch Einermengen]
\label{b28guide:thm:powerset}
$\mathcal P(S)$ und $\mathcal P_+(S)$ sind Halbgruppen.
Für beliebige Teilmengen und mengenindizierte Familien gelten
\[
 (AB)C=A(BC),\qquad
 A\left(\bigcup_i B_i\right)=\bigcup_i AB_i,\qquad
 \left(\bigcup_i A_i\right)B=\bigcup_i A_iB.
\]
Die Abbildung $a\mapsto\{a\}$ bettet $S$ als Unterhalbgruppe
in $\mathcal P(S)$ ein.
\cite{projectpower,howie}
\end{theorem}
